\section{Fazit und Ausblick} \label{sec:fazit}

\subsection{Zusammenfassung} \label{sec:zusammenfassung}

Der vorliegende Artikel hat gezeigt, dass das Gesundheitswesen bei der Einführung
KI-gestützter Bildanalyse auf etablierte Methoden und Systemarchitekturen der
industriellen Qualitätskontrolle zurückgreifen kann. Drei Kernaussagen lassen sich
festhalten:

\begin{enumerate}
    \item \textbf{Methodischer Transfer ist empirisch belegt.} CNN-Architekturen und
    Transfer Learning erzielen auf medizinischen Bilddaten radiologische Befundqualität;
    CheXNet übertrifft den Radiologen-Durchschnitt statistisch signifikant
    \cite{chexnet2017}.

    \item \textbf{Nutzen ist quantitativ messbar.} Der kombinierte Einsatz von KI und
    Radiologe reduziert den Befundungsaufwand um bis zu 44~\% \cite{breastScreening2025}
    und senkt das Fehlerrisiko bei erhöhter Arbeitslast \cite{workloadErrors2023}.

    \item \textbf{Datenschutzkonforme Architektur ist verfügbar.} Federated Learning mit
    Differential Privacy sowie On-Premise-Infrastrukturen erfüllen die Anforderungen von
    DSGVO und EU AI Act \cite{euAiAct2024, federatedLearning2022}.
\end{enumerate}

\subsection{Ausblick} \label{sec:ausblick}

Drei Entwicklungen werden die weitere Konvergenz beider Domänen prägen: Multimodale
Modelle, die Bilddaten mit klinischen Texten und Labordaten verbinden; die regulatorische
Reife durch den EU AI Act, der einheitliche Qualitätsstandards für diagnostische
KI-Systeme schafft \cite{euAiAct2024}; sowie Federated Learning als Standard für
datensouveränes, standortübergreifendes Modelltraining in Klinikverbünden
\cite{federatedAdvances2025}. Zusammenfassend zeigen die untersuchten Verfahren und
Architekturen, dass der Transfer industrieller KI-Bildanalyse in das Gesundheitswesen
methodisch fundiert und regulatorisch umsetzbar ist.
